% Towards Expert Financial QA via Self-Improving RAG
% FINAI@ICLR 2026 Workshop Paper (4-page version)

\documentclass{article}
\usepackage{iclr2026_conference,times}

% Required packages
\usepackage{hyperref}
\usepackage{url}
\usepackage{booktabs}
\usepackage{amsfonts}
\usepackage{amsmath}
\usepackage{amssymb}
\usepackage{amsthm}
\newtheorem{proposition}{Proposition}
\usepackage{bbm}  % For indicator function \mathbbm{1}
\DeclareMathOperator*{\argmax}{arg\,max}
\usepackage{nicefrac}
\usepackage{microtype}
\usepackage{graphicx}
\usepackage{xcolor}
\usepackage{tikz}
\usetikzlibrary{positioning, arrows.meta, shapes.geometric, fit}
\usepackage{algorithm}
\usepackage{algorithmic}
\usepackage{enumitem}
\usepackage[most]{tcolorbox}

% Title
\title{Towards Expert Financial QA via Self-Improving RAG}

% Authors
\author{
  Junjie Xiong \\
  University of California, Berkeley \\
  \And
  Shawheen Ghezavat \\
  California Polytechnic State University \\
  \And
  Aum Hirpara \\
  Hofstra University \\
  \And
  Sean Wu \\
  Pepperdine University
}

% Uncomment for camera-ready (accepted):
\iclrfinalcopy

\begin{document}
\raggedbottom  % Prevent excessive vertical spacing between sections

\maketitle

\begin{abstract}
Expert-level financial question answering requires both \textbf{grounded verification} to catch numeric hallucinations and \textbf{audit trails} for regulatory compliance, attributes that standard single-pass RAG systems lack. We take a step toward this goal with Self-Improving RAG, a framework that decomposes document QA into three specialized agents (Retrieval, Reasoning, and Judge) coordinated by an orchestrator with feedback-driven self-correction. When the Judge Agent scores an answer below a dynamic threshold, the system triggers retry with escalated strategies: broader retrieval, more careful prompting, and relaxed acceptance criteria. We evaluate on FinanceBench (SEC filing QA), where Self-Improving RAG achieves 86\% oracle-guided accuracy (measuring agreement with gold answers) with a 36.4\% Lazarus Rate, recovering nearly 4 in 10 initially incorrect answers through targeted retry. A key finding is that a fixed retrieval pipeline with judge-driven retry achieves strong results without dynamic routing, providing full interpretability. Every decision is logged with confidence scores, enabling the audit trails required for regulated financial applications.
\end{abstract}

% Main paper sections (4 pages excluding references)
\input{intro}
\input{related}
\input{method}
\input{results}
\input{conclusion}

% Bibliography
\bibliography{references}
\bibliographystyle{iclr2026_conference}

% Appendix (does not count toward page limit)
\appendix
\input{appendix}

% LLM Usage Disclosure
\section*{Statement on LLM Usage}

In accordance with ICLR's policy on Large Language Models, we declare:

\textbf{Text Refinement:} LLMs assisted with grammar and clarity improvements.

\textbf{Citation Verification:} We developed an automated citation verification
agent that cross-referenced all claims against source abstracts via Semantic
Scholar to ensure citation accuracy.

\textbf{Figure Design:} AI tools assisted with figure design and layout.

All technical contributions, experimental results, and analysis were conceived
and verified by the authors.

\end{document}
